\documentclass[12pt]{article}

%\textheight 9in
%\textwidth 6in
%\topmargin 0in
%\oddsidemargin 0.5in
%\evensidemargin 0.5in

\newcounter{exno}
\newcommand{\exnume}{\par\smallskip
\noindent\addtocounter{exno}{1}%
{\bf Esercizio \theexno.\ }}
\title{Esercizio Laboratorio 2}


\begin{document}
\maketitle

\noindent
Produrre un documento (con contenuto di senso compiuto, possibilmente matematico!) di max 2 pagine che contiene:

\begin{enumerate}
  \item del testo centrato / a destra /a sinistra
  \item del testo di almeno $3$ misure diverse
  \item elenco puntato o numerato
  \item una descrizione per punti
  \item una tabella con (semplici) formule matematiche
  \item una tabella articolata con multicolumn
  
\end{enumerate}

\noindent
\textbf{Esempi}: Relazione sulla scorsa lezione di CAM1 - Relazione sull'ultimo esame sostenuto - Relazione sul mio percorso di laurea triennale - Relazione sul mio percorso di laurea Magistrale - Relazione su statistiche varie del corso di laurea in matematica - Piccola guida per gli studenti che si devono iscrivere in matematica - Piccola lezione di matematica per dei bambini di prima media - ....


\end{document} 